\documentclass[11pt,a4paper]{article}
\usepackage[T1]{fontenc}
\usepackage[utf8]{inputenc}
\usepackage[english]{babel}
\usepackage{amsmath,amssymb}
\usepackage[margin=2cm]{geometry}
\setlength{\parindent}{0pt}
\setlength{\emergencystretch}{3em}
\allowdisplaybreaks

\title{About \texttt{growth\_scale\_limits.tex}}
\author{}
\date{}
\begin{document}
\maketitle
This note explains what the test data in \texttt{growth\_scale\_limits.tex} is used for. That file is input for the engine in \texttt{main.c}; it is read by the test \texttt{test\_limits.c} in this folder.

\begin{enumerate}
\item \textbf{Purpose.} The file holds 34 true statements in which exponentials, powers and logarithms compete: $\frac{e^{x}}{x^{5}}$, $x^{100} e^{-x}$, $\frac{(\ln x)^{100}}{x}$, $2^{x} - 3^{x}$, $\frac{x^{x}}{e^{x^{2}}}$, $(1 + e^{x})^{\frac{1}{x}}$, $\frac{\ln \ln x}{\sqrt{\ln x}}$, and beside $0$ the functions $e^{\pm\frac{1}{x}}$ against powers of $x$. They test the scale of growth of chapter 29 (rule 29k: a factor $e^{G}$ beside the expansion, compared factor by factor) and the substitution $x = e^{\frac{1}{t}}$ for the logarithms of logarithms (rule 29l), with no exercise of a book behind them: they are the classic examples of the hierarchy of infinities.

\item \textbf{Where the values come from.} Each value follows from the hierarchy (an exponential $e^{H}$ with $H \to +\infty$ faster than $\ln x$ beats every power, a power beats every power of a logarithm) and was computed by hand; for the quotients and the exponential forms the leading terms are written in the comment of rule 29k. The program does not know them in advance.

\item \textbf{Expected result.} The line \texttt{PASS: 34 of 34 statements proved, each with a written proof}, and the same statements asked as requests (without their value) answered with a limit that is then proved.

\item \textbf{How to run it.} From the top folder of the repository:
\begin{verbatim}
cc -std=c11 -O2 -Wall -Wextra -o /tmp/test_limits \
    tests/limits/test_limits.c && /tmp/test_limits
\end{verbatim}
\end{enumerate}
\end{document}
